% -----------------------------------------------------------------
\subsection{Functional Test Suite}
\label{subsec:crummy-test}

\code{CrummyTest} is a self-contained functional test suite invoked
from the application's right-click context menu
(\code{IDM\_BASICTEST}, labeled ``Run Basic Test'').  The suite
exercises the direct rendering path
(\cref{subsec:two-render-paths}): each test calls \code{Drain()}
first to ensure no pool workers remain active, then invokes
\code{CalcFractal(true)}, which writes results to \code{m\_CurIters}.

The top-level \code{TestAll()} method runs the full suite in sequence:
\begin{itemize}
  \item \code{TestWindowResize}---resizes the application window and
        verifies that a re-render completes correctly.
  \item \code{TestBasic}---iterates through every
        \code{RenderAlgorithm} entry, rendering the default view,
        View~5, and View~11 with both 32-bit and 64-bit iteration
        types.
  \item \code{TestReferenceSave}---saves and reloads reference orbits,
        verifying round-trip correctness.
  \item \code{TestVariedCompression}---tests reference-orbit
        compression at varying error thresholds.
  \item \code{TestImaginaLoad}---loads orbits stored in the Imagina
        format (\cref{subsec:imagina-format}).
  \item \code{TestStringConversion}---validates string-to-precision
        and precision-to-string conversion utilities.
  \item \code{TestPerturbedPerturb}---exercises perturbed iteration
        chains for numerical correctness.
  \item \code{TestGrowableVector}---validates the
        \code{GrowableVector} data structure
        (\cref{sec:disk_backed_growable_vectors}).
\end{itemize}
A separate \code{TestReallyHardView27} (not included in
\code{TestAll}) renders View~\#27, a period ${\sim}28$~billion point
requiring hours of computation; the test serves as an overnight
stress benchmark rather than a routine regression check.

% -----------------------------------------------------------------
\subsection{Kernel List}
\label{sec:kernel-list}

This section simply lists all the CUDA kernels in \FractalShark{}.

\begin{itemize}

\item \textbf{Mandelbrot base kernels} (\cref{sec:mandel-base-kernels})
  \begin{itemize}
    \item \code{mandel\_1x\_float}
    \item \code{mandel\_1x\_double}
    \item \code{mandel\_2x\_float}
    \item \code{mandel\_2x\_double}
    \item \code{mandel\_4x\_float}
    \item \code{mandel\_4x\_double}
    \item \code{mandel\_hdr\_float}
  \end{itemize}

\item \textbf{Mandelbrot perturbation kernels}
  \begin{itemize}
    \item \code{mandel\_1x\_double\_perturb\_bla}
    \item \code{mandel\_1xHDR\_float\_perturb\_bla} (\cref{sec:bilinear-approx})
    \item \code{mandel\_1xHDR\_float\_perturb\_lav2} (\cref{sec:la-v2-perturb,sec:perturb-only})
  \end{itemize}

\item \textbf{High-precision reference kernels} (\cref{subsec:ref-orbit-gpu})
  \begin{itemize}
    \item \code{HpSharkReferenceGpuKernel}
    \item \code{HpSharkReferenceGpuLoop}
  \end{itemize}

\item \textbf{Utility kernels}
  \begin{itemize}
    \item \code{antialiasing\_kernel}
    \item \code{max\_reduce}
    \item \code{max\_kernel}
  \end{itemize}

\item \textbf{Addition kernels} (\cref{sec:hp-add})
  \begin{itemize}
    \item \code{AddKernel}
    \item \code{AddKernelTestLoop}
  \end{itemize}

\item \textbf{Multiplication (NTT) kernels} (\cref{sec:ntt-multiply})
  \begin{itemize}
    \item \code{MultiplyKernelNTT}
    \item \code{MultiplyKernelNTTTestLoop}
  \end{itemize}

\item \textbf{Disabled kernels}
  \begin{itemize}
    \item \code{mandel\_2x\_float\_perturb\_setup}
    \item \code{mandel\_2x\_float\_perturb}
  \end{itemize}

\end{itemize}

\subsubsection{Commentary}

\begin{itemize}

\item All the per-pixel CUDA kernels are in \code{gpu\_render.cu}. The two
most interesting are probably \code{mandel\_\allowbreak 1xHDR\_\allowbreak float\_\allowbreak perturb\_bla} and
\code{mandel\_\allowbreak 1xHDR\_\allowbreak float\_\allowbreak perturb\_lav2}, which are the ones I've spent the
most time on lately. For better performance at low zoom levels, consider
\code{mandel\_1x\_float\_perturb}, which leaves out linear approximation
and just does straight perturbation up to $\sim 10^{30}$, which corresponds
with the 32-bit float exponent range.

\item The \code{mandel\_1x\_float} (\cref{sec:mandel-1x-float}) is the classic 32-bit float Mandelbrot and 
is screaming fast on a GPU. This one is optimized with fused multiply-add for 
fun even though it is of limited practical use because one can barely zoom in before
pixellation sets in.

\end{itemize}

